% !TEX root = ./main.tex
\chapter{Valutazione tecnica del prototipo}
\label{chap:risultati}

Dopo aver descritto il prototipo nei capitoli precedenti, qui interessa capire come si comporti davvero nelle prove dal punto di vista tecnico. L'attenzione è sulla tenuta della pipeline conversazionale, sui tempi osservati, sui failure case e sui limiti che emergono quando il sistema passa dal caso nominale all'uso reale. La parte comparativa sugli utenti viene invece ripresa nel capitolo successivo.

\section{Impostazione della valutazione}
\label{sec:metodologia}
\noindent Il primo passo era fissare che cosa osservare davvero e in quali condizioni farlo. Per questo la valutazione è stata costruita attorno agli scenari di prova, ai punti di misura più utili per leggere il turno conversazionale e a una batteria tecnica separata sul backend, usata per isolare meglio il comportamento del RAG. Il quadro architetturale è discusso in \autoref{sec:architettura-soulframe}; qui contano soprattutto le osservazioni raccolte usando il prototipo.

Le prove sono state svolte in due ambienti funzionalmente equivalenti. In locale, il client Unity in \texttt{MainMode} comunica con i micro-servizi su loopback (\texttt{127.0.0.1:8001--8004}). Nell'ambiente Ubuntu usato per le prove, la stessa pipeline è esposta dietro reverse proxy con path \texttt{/api/*}. La logica applicativa è la stessa; cambia il contesto operativo.

L'ambiente locale Windows è un laptop con CPU \texttt{AMD Ryzen AI 9 HX 370} e dGPU \texttt{NVIDIA RTX 5070 Ti Mobile} con 12 GB di VRAM. L'ambiente Ubuntu di prova gira su Google Cloud, con GPU \texttt{NVIDIA L4} (VRAM tipica 24 GB), disco persistente da 60 GB, IPv4 statico e dominio dedicato ai test. In questa configurazione, lo STT usa Whisper \texttt{small} come default in locale Windows e Whisper \texttt{medium} sul server Ubuntu, seguendo la diversa disponibilità di VRAM.

Il flusso osservato segue il percorso operativo del prototipo: \texttt{POST /transcribe} per la trascrizione, \texttt{POST /chat} per retrieval e generazione, \texttt{POST /tts\_stream} per la risposta vocale in streaming. Il confronto tra i due ambienti è stato condotto su sessioni omogenee per avatar, lingua e stato memoria (presente/assente), in modo da non far dipendere le differenze da configurazioni troppo diverse.

I punti di osservazione sono stati fissati in modo uniforme: fine registrazione, completamento STT, completamento RAG, primo audio TTS e inizio playback. I tempi non derivano da una telemetria automatica per turno, ma da log di servizio, stato UI ed esiti HTTP raccolti su un numero limitato di sessioni esplorative. Servono soprattutto a capire l'ordine di grandezza dei colli di bottiglia e a confrontare le due configurazioni.

Gli scenari conversazionali coprono sia il percorso nominale sia i casi più sensibili alla latenza percepita: domanda aperta, interrogazione di memoria per-avatar, sequenze brevi e ravvicinate di turni, sessioni prolungate. Prima di ogni sessione è stata verificata la disponibilità dei servizi tramite endpoint \texttt{/health}, per evitare che il confronto fosse alterato da boot incompleti o configurazioni incoerenti.

Per leggere meglio il comportamento del RAG è stata affiancata anche una verifica tecnica separata sul solo backend locale Windows. In questo caso il client Unity, lo STT e il TTS non entrano nella prova: il servizio viene interrogato via HTTP con profili sintetici costruiti per stressare memoria testuale, immagini, documenti, collisioni di sorgente e coerenza identitaria.

Le osservazioni derivano dalle risposte JSON, dallo stato del servizio e dai contenuti effettivamente recuperati; i log conversazionali per avatar/sessione hanno aiutato a rileggere alcuni casi, ma non sostituiscono una telemetria turn-based completa e normalizzata. Una sintesi dei risultati tecnici più recenti è riportata in Appendice~\ref{app:notes}.

\subsection{Metriche e criteri di valutazione}
La valutazione tecnica guarda soprattutto a reattività della pipeline e tenuta del turno conversazionale. Le misure usate sono: tempo STT, tempo del blocco RAG/LLM, tempo al primo audio della fase TTS, latenza end-to-end del turno e \emph{cold-start/time-to-ready} dello stack. In questa fase non è stata usata una strumentazione persistente turn-by-turn. I tempi STT, RAG/LLM, TTS ed end-to-end derivano da osservazione controllata e lettura dei log. Il dato più solido resta il tempo al primo audio TTS; per RAG/LLM e per il tempo complessivo del turno i valori vanno letti come stime coerenti con i test esplorativi. {\setlength{\emergencystretch}{5em}\hbadness=10000\sloppy Per la robustezza operativa si osservano richieste fallite, timeout lato client e necessità di ripetere il turno. Con la policy attuale, che attende circa quindici secondi e prova un solo recupero, questa misura dà soprattutto un'idea di quanto spesso la pipeline si interrompa o degradi durante l'uso continuativo.\par}

Per la qualità STT, la Word Error Rate (WER) è ancora la metrica di riferimento in letteratura, ma qui non è stata misurata su un corpus controllato: le osservazioni riguardano soprattutto il comportamento di Whisper \texttt{small} in locale Windows e \texttt{medium} nell'ambiente Ubuntu di prova. Per il RAG distinguo invece due aspetti: pertinenza del contesto recuperato e fedeltà della risposta rispetto alle fonti che il backend ha effettivamente usato; quando serve, questo controllo può essere ricostruito anche dall'elenco delle sorgenti recuperate riportato nei log \cite{gao2024ragsurvey,Radford2023Whisper}.

Accanto alle sessioni end-to-end, sono state usate due batterie tecniche sul solo backend. La prima, multi-sorgente, serviva a capire se il modello piccolo riuscisse a tenere insieme memoria testuale, immagini e documenti senza perdere la sorgente giusta. La seconda, solo testuale, serviva invece a verificare un punto più delicato: se la risposta finale restasse coerente con l'identità dell'avatar senza confondere il profilo del personaggio con quello dell'utente.

Nelle prove miste, che combinano memoria testuale, immagini e documenti, il retrieval ha retto in tutti i casi osservati (28 su 28). Il \emph{full pass} si è fermato a 24 casi su 28: in altre parole, il sistema trova quasi sempre il materiale corretto, ma in quattro casi non riesce ancora a trasformarlo in una risposta abbastanza precisa. Un caso riuscito è la domanda sul voto di Fisica 1, in cui il PDF corretto viene recuperato e la risposta restituisce il 30/30. Un caso più debole è invece la richiesta di elencare tutti gli esami presenti nel documento: il file viene trovato, ma il riepilogo finale resta troppo generico. Nella batteria solo testuale dedicata alla coerenza identitaria, retrieval e correttezza fattuale hanno tenuto in 22 casi su 22, con esito complessivo pari a 19 casi su 22. Qui i cedimenti principali non riguardano il fatto recuperato, ma la riformulazione finale: in una prova l'avatar Viola si presenta correttamente, ma nella stessa risposta ingloba anche un dettaglio che appartiene all'utente. Il dettaglio per categorie e casi limite è raccolto in Appendice~\ref{app:notes}.

Accanto alla dimensione tecnica, la valutazione include anche una componente di esperienza utente. Gli indicatori UX richiamano il quadro discusso nel Capitolo~\ref{chap:stato-arte}, soprattutto \autoref{sec:eca-xr} e \autoref{sec:pipeline-ai}. Contano tre aspetti: chiarezza del turn-taking, comprensibilità del parlato sintetizzato e leggibilità dello stato conversazionale in interfaccia. Le osservazioni qualitative seguono indicatori espliciti: frequenza di ripetizioni o riformulazioni, numero di turni di riparazione, tolleranza all'attesa, stabilità del ritmo conversazionale, coerenza percepita tra voce e avatar. Per il sottosistema visivo considero anche chiarezza degli stati (listening/processing/speaking), comfort visivo (flicker, intensità bloom, movimenti camera) e coerenza stilistica complessiva. Il target di valutazione resta ampio e non è ristretto a una singola fascia anagrafica. Gli accorgimenti di accessibilità (icone dinamiche, feedback testuali e segnali visivi persistenti) vengono quindi letti come requisiti di leggibilità trasversali, non come adattamenti pensati per un solo profilo utente.

\section{Risultati tecnici del prototipo}
\label{sec:quantitativi}
\noindent I risultati tecnici hanno senso soprattutto se letti su due piani distinti. Da una parte c'è il costo della pipeline STT--RAG--TTS nel turno tipico; dall'altra c'è il modo in cui quel costo cambia quando entrano in gioco avvio dei servizi, rete, proxy e differenze tra build Windows e setup Ubuntu. Le prove tecniche mostrano che la pipeline STT--RAG--TTS è utilizzabile in entrambi gli ambienti testati. La tabella seguente condensa l'ordine di grandezza dei tempi osservati, mettendo in evidenza dove si concentra il costo della pipeline. I valori sono orientativi: derivano da osservazioni esplorative e log di esecuzione, non da una campagna strumentata per turno. Le scelte implementative di riferimento sono nel Capitolo~\ref{chap:implementazione}, soprattutto per i moduli RAG e TTS (\autoref{subsec:rag-memory}, \autoref{subsec:tts-streaming}).

\par\begingroup
\small
\setlength{\LTleft}{0pt}
\setlength{\LTright}{0pt}
\begin{longtable}{|>{\raggedright\arraybackslash}p{0.28\textwidth}|>{\raggedright\arraybackslash}p{0.33\textwidth}|>{\raggedright\arraybackslash}p{\dimexpr0.39\textwidth-4\tabcolsep-3\arrayrulewidth\relax}|}
\caption[Tempi per fase della pipeline]{Tempi indicativi per fase della pipeline.}
\label{tab:tempi-fase} \\
\hline
\textbf{Fase} & \textbf{Locale (Windows)} & \textbf{Setup Ubuntu testato} \\
\hline
\endfirsthead
\hline
\textbf{Fase} & \textbf{Locale (Windows)} & \textbf{Setup Ubuntu testato} \\
\hline
\endhead
STT (Whisper: \texttt{small}/\texttt{medium}) & nell'ordine di 1--3 s & nell'ordine di 1--4 s \\
\hline
RAG+LLM (Ollama) & circa 3 s nella maggior parte dei turni osservati & circa 3 s nella maggior parte dei turni osservati \\
\hline
TTS streaming (tempo al primo audio) & circa 5 s (tempo al primo audio) & circa 7--8 s (tempo al primo audio) \\
\hline
End-to-end (rilascio PTT $\rightarrow$ primo audio) & circa 9--12 s nei turni tipici osservati & circa 11--15 s nei turni tipici osservati \\
\hline
\end{longtable}
\par\endgroup

La fase STT è risultata sensibile soprattutto alla durata dell'audio in ingresso, con intervalli osservati di 1--3 s in locale Windows e 1--4 s sul server Ubuntu di prova. Il blocco RAG/LLM, pur senza una strumentazione sistematica per turno, si è collocato spesso attorno ai 3 s in entrambi gli ambienti, con oscillazioni legate soprattutto al tipo di query e al carico del runtime locale di inferenza. Il collo di bottiglia percepito è comunque la fase TTS in termini di tempo al primo audio: circa 5 s in locale e circa 7--8 s sul server Ubuntu di prova. Il tempo end-to-end del turno si muove di conseguenza intorno a 9--12 s in locale e 11--15 s sul server Ubuntu di prova nei casi tipici osservati, ma questo dato va considerato orientativo perché non nasce da una telemetria automatica persistente.

\subsection{Latenza, stabilità e differenze ambientali}
Cold-start e turni successivi vanno tenuti distinti: il primo dice quanto tempo serve prima di poter iniziare, il secondo dice quanto il sistema si mantiene regolare una volta avviato. Le criticità osservate sono coerenti con quanto discusso in \autoref{subsec:latenza-timeout} e \autoref{subsec:error-retry-fallback}.

\begin{table}[t]
\centering
\small
\begin{tabularx}{\linewidth}{l>{\centering\arraybackslash}X>{\centering\arraybackslash}X}
\hline
\textbf{Misura} & \textbf{Locale (Windows)} & \textbf{Setup Ubuntu testato} \\
\hline
Pre-attivazione/setup iniziale (stack pronto) & circa 40 s & circa 80 s \\
\hline
\end{tabularx}
\normalsize
\caption[Tempo di inizializzazione]{Tempo di inizializzazione osservato.}
\label{tab:cold-start-ready}
\end{table}

I valori in Tabella~\ref{tab:cold-start-ready} descrivono solo la fase iniziale di pre-attivazione/setup. Nel setup osservato, l'ordine di grandezza è stato di circa 35 s per la pre-attivazione di Coqui XTTS v2 e circa 5 s per la pre-attivazione RAG in locale Windows, contro circa 60 s e circa 20 s nell'ambiente Ubuntu usato per le prove. Sono stime ricavate dall'avvio dei servizi, utili per leggere dove si concentra l'attesa iniziale più che per scomporre in modo rigoroso ogni turno.

Dopo la pre-attivazione iniziale il peso si sposta altrove: nei turni successivi contano soprattutto tempo al primo audio TTS, rete e scambio dati tra client e servizi. Per ottenere una scomposizione più solida della latenza servirebbe una telemetria turn-based persistente, con ID turno e timestamp sincronizzati tra client e servizi. Nello stato attuale i log bastano per un confronto attendibile tra ambienti, ma non per una decomposizione fine di ogni singolo turno.

Sul fronte stabilità, i problemi più ricorrenti nel campione osservato sono timeout client-side nei turni più lenti e degradazione prestazionale quando le risorse si riducono. Nelle sessioni esaminate non sono emersi blocchi strutturali della macchina a stati del client. Nell'ambiente Ubuntu di prova, però, i tempi oscillano di più perché a inferenza e TTS si sommano overhead di rete e proxy.

Il confronto tra ambiente locale e server Ubuntu di prova mostra che il contratto applicativo non cambia, ma il comportamento operativo sì. In locale quasi tutto avviene sulla stessa macchina, quindi molte richieste stanno nell'ordine di pochi millisecondi e il tempo percepito riflette soprattutto l'inferenza. Sul server Ubuntu si sommano invece upload e download reali: lo si nota nell'upload del file audio verso STT, nel download dell'avatar quando non è ancora in cache, nello stream vocale progressivo e anche nelle altre richieste HTTP del flusso. Un avatar Avaturn pesa circa 16 MB; quando gli asset sono già presenti in cache una parte del costo sparisce e il turno si alleggerisce un poco. Nei casi peggiori, invece, la somma di rete, proxy e inferenza rende più visibili jitter e oscillazioni tra un turno e l'altro.

In questo contesto è emersa anche una criticità più specifica sul sottosistema vocale: oltre alla latenza, il playback TTS WebGL può introdurre artefatti intermittenti che in locale risultano assenti o molto meno evidenti. Il problema non coincide solo con la qualità della sintesi Coqui, ma con il comportamento dell'intera catena di riproduzione progressiva nel browser remoto, dove proxy, jitter e vincoli del runtime audio WebGL diventano parte del problema percepito. Per la stessa ragione, nella configurazione attuale il \emph{word-by-word reply flow} non viene più mantenuto attivo in WebGL: il client usa un fallback testuale statico e mostra subito il reply completo, per evitare altra instabilità.

Nel confronto diretto con la build Windows dello stesso client la differenza è emersa in modo abbastanza netto. Il lag audio percepito che in WebGL tende a comparire nei turni più sensibili non si è ripresentato con la stessa evidenza. A parità di backend, pipeline STT--RAG--TTS e logica applicativa, cambia soprattutto il contesto di esecuzione del playback. In WebGL il browser aggiunge una mediazione in più, tra buffering, scheduling e tolleranza al jitter. In Windows la riproduzione resta più lineare e il turno suona più continuo. La stessa build Unity mantiene coerenza funzionale nei due contesti grazie alla normalizzazione degli endpoint; nel deploy pubblico diventano però più evidenti health-check, readiness reale dei servizi e controllo dei colli di bottiglia inferenziali.

\section{Risultati qualitativi e casi d'uso}
\label{sec:qualitativi}
\noindent Accanto ai tempi e ai failure più evidenti, serviva anche capire come il prototipo venisse percepito quando il turno riusciva davvero a partire e a chiudersi. In questa parte la lettura si sposta quindi su qualità della voce, leggibilità dell'interfaccia e casi concreti in cui retrieval, risposta finale e feedback di stato restano allineati oppure iniziano a divergere. La qualità percepita dell'interazione vocale dipende soprattutto da tre fattori: naturalezza, reattività e continuità del turno. La letteratura sulla \emph{Voice User Experience} misura queste dimensioni in modi molto diversi tra loro; qui contano soprattutto le osservazioni raccolte sul prototipo.\cite{seaborn2021voiceux} I dettagli implementativi sono nel Capitolo~\ref{chap:implementazione} (\autoref{subsec:mainmode}, \autoref{subsec:tts-streaming}, \autoref{subsec:ps2-postprocessing}).

Per la voce sintetizzata, l'impiego di Coqui XTTS v2 con \emph{voice cloning} rafforza subito la percezione di identità dell'avatar. Il preprocessing del testo (\texttt{\_clean\_tts\_text}) evita letture troppo letterali di simboli e punteggiatura, mentre i parametri di campionamento e penalità di ripetizione sono tarati per privilegiare stabilità e scorrevolezza. La somiglianza con il campione originale non è costante, ma già così la voce risulta abbastanza riconoscibile da sostenere l'identità del personaggio durante il turno.

Quando la memoria è attiva, la coerenza delle risposte si nota soprattutto nella capacità dell'avatar di ``parlare come la persona''. Il servizio RAG distingue contenuti fattuali e tracce di stile, ricostruisce note di stile del personaggio e costruisce un riepilogo compatto di indizi utili da affiancare al contesto recuperato. Il prompt di sistema impone anche risposte in prima persona ed evita etichette meta o indicazioni sceniche. Il vincolo di \emph{grounding}, cioè l'obbligo di restare aderenti ai contenuti davvero recuperati, insieme alla possibilità di rispondere ``Non ricordo'' quando il contesto non basta, rende l'agente meno brillante ma più credibile: in un contesto embodied, un'incertezza esplicita pesa meno di un dettaglio inventato.

La batteria RAG eseguita separatamente sul backend aiuta a distinguere meglio retrieval diretto e conversazione. Con \texttt{/recall}, i ricordi manuali espliciti e le parafrasi semplici risultano in genere recuperabili in modo abbastanza stabile: preferenze, dettagli biografici e informazioni rare riemergono tra i primi candidati con buona coerenza rispetto all'avatar interrogato. Il comportamento si fa più incerto quando la query è vaga o semanticamente vicina a un fatto che appartiene all'altro avatar: in quei casi l'isolamento per-avatar regge spesso, ma non in modo assoluto. Il confronto tra \texttt{/recall} e \texttt{/chat} evidenzia anche un altro punto: una base di memoria che in retrieval diretto appare leggibile può produrre, in chat, una risposta finale meno solida. Alcune query fattuali dirette vengono instradate come \texttt{chitchat}, con perdita del contesto RAG disponibile; in altri casi la risposta conclusiva introduce dettagli non ben sostenuti dai ricordi recuperati. Il limite, quindi, non sembra stare solo nel retrieval, ma anche nel routing dell'intent e nella riscrittura finale.

La fluidità del turno dipende anche da come il sistema gestisce i tempi morti tra fine input e inizio output vocale. In \texttt{MainMode} il client aggiorna lo stato con messaggi come ``Sto pensando...'' durante la richiesta a \texttt{/chat} e ``Sto parlando...'' quando la risposta entra nella fase vocale. Il reply viene poi ricevuto tramite \texttt{/tts\_stream} come stream WAV progressivo, così la riproduzione non è costretta ad attendere l'intero trasferimento. Per ridurre il silenzio percepito nella finestra iniziale, il client riproduce una \emph{wait phrase} breve, selezionata da un insieme predefinito di chiavi e scaricata on-demand dall'endpoint \texttt{/wait\_phrase}. Quando il reply vero è pronto a partire, quel segnale di attesa viene interrotto e lascia spazio alla voce dell'avatar. La stessa logica supporta l'interruzione: quando l'utente riprende la parola, il client può abortire richieste e playback in corso evitando sovrapposizioni tra audio tardivo e nuovo turno. La diagnostica resta disponibile tramite \texttt{DebugText} nella \texttt{Status Bar} (disattivabile con tasto \texttt{INS} su desktop o con triplo tap della barra in modalità touch), separando il feedback utente dai dettagli di debug. Il sottosistema visivo lavora soprattutto su atmosfera e continuità percettiva, non come canale diagnostico principale. Lo stato del turno viene comunicato prima da testo di stato, transcript/reply e \texttt{UIHintBar}; rings e testo azzurro lampeggiante (desktop) accompagnano questi segnali senza sostituirli.\cite{li2021pueva}

Nella stessa batteria mirata, i materiali non testuali sono stati utili per capire dove il retrieval tiene e dove invece si indebolisce. Una richiesta come ``cosa ricordi del diagramma client server?'' è stata gestita bene sia in \texttt{/recall} sia in \texttt{/chat}: il sistema ha richiamato la descrizione dell'immagine e l'ha usata per produrre una risposta centrata. Con i PDF il quadro è più fragile. La domanda ``quali esami risultano nel documento?'' recupera il file corretto, ma la risposta finale tende a comprimere troppo il contenuto e a restare vaga. Quando l'OCR produce chunk pochi e ordinati il documento resta interrogabile; quando invece cresce il rumore OCR, il contesto recuperato si indebolisce e con esso anche la risposta.

\begin{figure}[H]
\centering
\includegraphics[width=0.95\textwidth]{mainmode_desktop_eval.png}
\caption[MainMode desktop]{Interfaccia desktop in modalità MainMode con transcript, risposta dell'avatar, stato conversazionale e comandi contestuali nella hint bar.}
\label{fig:mainmode-desktop}
\end{figure}

Figura~\ref{fig:mainmode-desktop} documenta un esempio di interazione in \texttt{MainMode} su desktop: la schermata mostra insieme trascrizione dell'utente, risposta dell'avatar, testo di stato e grammatica dei comandi nella hint bar.

L'usabilità dell'interfaccia è stata osservata confrontando due varianti che condividono lo stesso flusso logico, ma cambiano molto nel modo di farsi usare. In entrambe le modalità, il comportamento dell'applicazione è governato dagli stessi stati (\texttt{Boot}, \texttt{MainMenu}, \texttt{AvatarLibrary}, \texttt{SetupVoice}, \texttt{SetupMemory}, \texttt{MainMode}) e attiva la stessa pipeline backend. Questo confronto riguarda la superficie di interazione, non il runtime: desktop e touch descrivono il modo d'uso, mentre Windows nativo e WebGL descrivono il contesto di esecuzione del client. La differenza percepita nasce quindi da due assi distinti; in questa sottosezione conta soprattutto il primo, cioè tastiera/mouse contro gesto touch.

\begin{wrapfigure}{l}{0.34\textwidth}
\vspace{-6pt}
\centering
\includegraphics[width=0.95\linewidth]{mainmode_touch_eval.png}
\captionsetup{singlelinecheck=false,justification=raggedright}
\caption[MainMode touch]{Interfaccia touch in modalità MainMode: Push-to-Talk, swipe tra transcript e reply e adattamenti per l'usabilità su schermo touch.}
\label{fig:mainmode-touch}
\vspace{-8pt}
\end{wrapfigure}

Nella variante desktop l'azione principale è point-and-click, con supporto costante della hint bar per comandi e shortcut. Questo assetto si presta a un uso prolungato perché sfrutta abitudini già consolidate: l'utente individua rapidamente il focus, mantiene controllo fine con il mouse e accelera le azioni ricorrenti con la tastiera, soprattutto nei passaggi di setup. Quando il mouse resta inattivo per qualche secondo, parte dell'interfaccia si attenua e lascia più spazio ai contenuti principali.

La variante touch riorienta invece l'interazione attorno a un gesto centrale facile da memorizzare: il Push-to-Talk in \texttt{MainMode}. Il pulsante dedicato mostra subito se il sistema è in ascolto o in attesa, riducendo la necessità di ricordare scorciatoie o interpretare indicazioni testuali, cosa importante su schermi piccoli e in contesti d'uso intermittenti. Su touch tenere transcript e risposta nella stessa vista risultava però affollato; per questo i due contenuti sono stati separati con uno swipe abbastanza netto da evitare cambi accidentali, accompagnato da una transizione breve ma leggibile. Quando compare input testuale, l'avatar viene anche attenuato leggermente, così il volto non compete con l'area di scrittura ma continua a restare presente nella scena. Anche il ridimensionamento dell'interfaccia va nella stessa direzione: mantenere leggibilità su schermi e densità di pixel diverse. In entrambe le varianti, il volto dell'avatar funziona come anchor percettivo; su desktop lo spazio disponibile permette di mostrare insieme transcript, reply e hint operativi, mentre su touch la stessa simultaneità rischia di saturare lo schermo. Desktop e touch mostrano quindi compromessi netti: desktop favorisce rapidità e trasparenza grazie a spazio informativo e hint persistenti; touch privilegia immediatezza del gesto, ma con minore densità simultanea di informazioni, come mostra la Figura~\ref{fig:mainmode-touch}.


\subsection{Analisi di casi e failure cases}
I meccanismi tecnici richiamati nei tre casi sono analizzati in dettaglio nel Capitolo~\ref{chap:implementazione}, soprattutto \autoref{subsec:rag-memory}, \autoref{subsec:tts-streaming} e \autoref{subsec:error-retry-fallback}. Nel corso della valutazione esplorativa sono tornate più volte tre situazioni utili per capire dove la pipeline regge e dove invece comincia a perdere coerenza. I casi descritti di seguito non riportano trascrizioni o dialoghi letterali, ma sintetizzano input, comportamento e cause tecniche osservate.

Il primo caso riguarda una domanda che richiede davvero la memoria per-avatar e in cui il contesto recuperato è pertinente e non ambiguo. Lo scenario tipico è una richiesta su un'informazione già ingestita, formulata in modo abbastanza specifico da attivare un retrieval mirato. In questa condizione \texttt{rag\_server.py} costruisce un prompt su chunk coerenti tra loro e l'output dell'LLM resta allineato sia ai fatti riportati nella memoria sia allo stile dell'avatar. Contribuisce anche la sanitizzazione finale della risposta (\texttt{\_sanitize\_chat\_answer}), che rimuove prefissi di parlante e stage directions evitando che il TTS vocalizzi annotazioni non destinate all'utente. Per l'utente il turno scorre bene e la risposta suona abbastanza naturale, oltre che attribuibile all'identità dell'avatar.

Il secondo caso compare quando l'informazione richiesta è pertinente alla memoria ma il retrieval produce un contesto rumoroso o parzialmente ridondante. Succede soprattutto con domande formulate in modo ampio, che attivano chunk simili tra loro ma non abbastanza selettivi. In queste condizioni la deduplicazione riduce parte del rumore, ma non sempre basta: la risposta rimane nel tema, però diventa meno precisa, più generica o leggermente incoerente nello stile. Qui la qualità del risultato dipende più dalla qualità del set di chunk selezionati che dal solo vincolo di grounding. Dal lato dell'esperienza, i feedback visivi continuano a segnalare correttamente lo stato della pipeline, ma non comunicano la qualità interna del contesto; l'utente percepisce quindi un turno formalmente riuscito ma meno centrato.

{\setlength{\emergencystretch}{2em}\hbadness=10000\hfuzz=20pt
Il terzo caso riguarda modalità di guasto legate a timeout e recupero incompleto, con impatto diretto sulla continuità dell'interazione. Lo scenario più semplice è un fallimento in STT: \texttt{whisper\_\allowbreak server.py} non implementa retry interno e, se \texttt{/transcribe} non risponde in tempo o fallisce per errori temporanei, l'onere ricade sul client Unity. Nelle prove è emersa anche una fragilità più semplice: registrazioni quasi vuote, silenziose o troppo brevi possono generare una trascrizione spuria pur senza un contenuto vocale utile. Per limitarne gli effetti, il frontend blocca i WAV sotto circa \SI{0.55}{\second} e, dopo \texttt{/transcribe}, scarta alcune stringhe note come non valide; resta però un presidio euristico lato client, utile a ridurre falsi positivi evidenti ma non a esaurire il problema. Con la policy attuale, che attende circa quindici secondi e prova un solo ritentativo lato client, un singolo recupero può non bastare in condizioni di carico elevato o con audio più lungo del previsto; il comportamento osservato è quindi una transizione verso stato di errore o un ritorno in listening senza trascrizione valida. In modo analogo, la fase TTS può degradare quando si verifica un errore GPU: \texttt{coqui\_tts\_server.py} tenta un fallback CUDA$\rightarrow$CPU, ma se la generazione non riesce a produrre una risposta vocale utilizzabile l'esperienza del turno si interrompe o si appoggia a un testo non sincronizzato con il parlato. Le wait phrases aiutano a rendere più leggibile l'attesa breve; in un failure prolungato, però, non bastano da sole a distinguere una normale lentezza da un guasto. In questi casi il comportamento dei feedback visivi diventa cruciale: se non differenzia abbastanza processing normale e processing in errore, l'utente non distingue tra lentezza e fault.\par}

\vspace*{\fill}
Figura~\ref{fig:failure-sequence} propone un diagramma di sequenza che sintetizza un failure case centrato su timeout in STT, ritentativo lato client ed eventuale uscita verso messaggio di errore o ritorno a listening, rendendo esplicite le dipendenze tra timeout, retry e feedback percettivo (stato UI e rings).

\begin{figure}[H]
\centering
\includegraphics[width=\textwidth]{failure_sequence.png}
\caption[Failure case con timeout e retry]{Diagramma di sequenza di un failure case con timeout, retry e ricadute visibili sul feedback di stato.}
\label{fig:failure-sequence}
\end{figure}

\begingroup
\setlength{\emergencystretch}{4em}
\hbadness=10000
\sloppy
\section{Discussione dei risultati}
\label{sec:discussione}
\noindent A questo livello conta leggere insieme ciò che ha retto meglio e ciò che continua a mostrare fragilità. La discussione non aggiunge nuove prove, ma prova a tenere nello stesso quadro modularità del prototipo, qualità del retrieval, limiti della voce sintetica e differenze operative tra i due ambienti osservati.

\paragraph{Punti di forza.}
La modularità è uno dei punti più solidi del prototipo: i servizi Speech-to-Text (STT), Retrieval-Augmented Generation (RAG) e Text-to-Speech (TTS) sono separati, raggiungibili tramite endpoint dedicati e verificabili con controlli di stato (\texttt{/health}). Questa scomposizione rende più semplice localizzare i colli di bottiglia e ha retto bene anche nel passaggio tra ambiente locale e setup Ubuntu testato. Dal lato esperienza, lo streaming TTS rende l'attesa più leggibile anche quando la sintesi richiede più tempo e permette al client di avviare il playback senza attendere l'intero trasferimento del reply. Le wait phrases non riducono il tempo di inferenza, ma evitano che il silenzio venga percepito subito come blocco del turno.

La qualità del recupero beneficia della ricerca ibrida (segnali lessicali BM25 + similarità vettoriale), che migliora la pertinenza in presenza di nomi propri, parole chiave e parafrasi. L'isolamento della memoria per-avatar e la deduplicazione dei chunk limitano interferenze e ridondanze. Anche i fallback aiutano a mantenere il ciclo utilizzabile quando le condizioni non sono ideali. La verifica esplorativa del solo backend RAG va nella stessa direzione: il retrieval diretto ha recuperato bene ricordi manuali, preferenze e descrizioni immagine anche quando la query non ripeteva esattamente le parole della memoria salvata. Nello stesso perimetro è emerso anche un caso di memoria automatica efficace: un'informazione affidata esplicitamente al sistema via \texttt{/chat} è stata salvata dal backend in modo automatico e poi recuperata correttamente in una richiesta successiva. Anche l'interfaccia contribuisce alla robustezza complessiva: desktop e touch condividono la stessa macchina a stati, ma usano affordance diverse in base al dispositivo. Questo riduce divergenze implementative e mantiene abbastanza stabile il modello mentale dell'utente.

\paragraph{Criticità osservate e limiti emersi.}
\label{subsec:criticita-osservate}

Durante le sessioni esplorative sono emersi tre problemi abbastanza ricorrenti. Il primo riguarda l'aderenza alla memoria. In alcuni turni l'avatar ha prodotto affermazioni non presenti nella memoria indicizzata, soprattutto con query vaghe e contesto recuperato debole. Il risultato è una risposta plausibile ma non allineata a quanto realmente memorizzato. Nella batteria RAG questo limite si è visto soprattutto nel passaggio da \texttt{/recall} a \texttt{/chat}: query fattuali dirette che in retrieval producevano candidati pertinenti sono state talvolta instradate come \texttt{chitchat}, e in altri casi la risposta finale si è allontanata dal contesto recuperato pur con fonti leggibili nei log.

Un secondo problema è emerso con i PDF più complessi. Con tabelle o layout multi-colonna, il retrieval è apparso meno affidabile rispetto a note lineari o testo semplice. Il sintomo ricorrente è la restituzione di frammenti parziali o poco ordinati anche con documento indicizzato. Più sfumato, ma comunque presente, è il tema dell'identità del profilo. Anche l'isolamento per-avatar va letto con prudenza: quando le memorie dei due profili contengono dettagli volutamente vicini, il sistema tende spesso a mantenere la separazione corretta; con query semanticamente prossime ma riferite all'altro avatar, però, può emergere il dettaglio più simile del profilo corrente invece di una dichiarazione esplicita di assenza. La voce sintetizzata è risultata comprensibile e riconoscibile come voce dell'avatar, ma con somiglianza variabile rispetto al campione originale. L'effetto non blocca l'interazione, però riduce la coerenza d'identità percepita nei turni più lunghi.

Dietro queste criticità restano limiti strutturali coerenti con i problemi di implementazione già tracciati nel Capitolo~\ref{chap:implementazione} (\autoref{subsec:latenza-timeout}, \autoref{subsec:ocr-ingestione}, \autoref{subsec:compatibilita-dipendenze}). Il limite più evidente è ancora il tempo al primo audio TTS (Tabella~\ref{tab:tempi-fase}). La fase RAG/LLM appare più compatta nei turni tipici, ma è anche quella meno isolabile con precisione perché dipende dal runtime locale di inferenza, dal profilo del modello e dalle risorse disponibili. Nell'ambiente Ubuntu usato per le prove si sommano inoltre overhead di rete e proxy, con jitter più marcato nei casi lenti e maggiore probabilità di timeout lato client. C'è poi un limite specifico della piattaforma WebGL. La consegna progressiva del WAV migliora il tempo percepito al primo audio, ma non offre lo stesso controllo sul playback continuo disponibile nella build Windows nativa provata sullo stesso impianto applicativo. In presenza di jitter reale, la risposta può quindi risultare meno uniforme di quanto accada nella build Windows dello stesso frontend: il sottosistema TTS resta utilizzabile, ma più fragile in WebGL, soprattutto nel deploy remoto, che nella configurazione Windows locale su loopback.

La valutazione dell'esperienza resta esplorativa: le conclusioni su naturalezza, co-presenza e impatto dei feedback visivi derivano da osservazioni qualitative. Sono utili per orientare le iterazioni progettuali, ma non bastano a sostenere inferenze robuste su efficacia o gradimento.\cite{seaborn2021voiceux}

Sul lato tecnico i limiti sono più facili da circoscrivere. La Word Error Rate (WER) non è stata misurata in modo sistematico su un corpus di riferimento. In locale Windows lo STT usa Whisper \texttt{small}. Nell'ambiente Ubuntu di prova usa \texttt{medium}, reso praticabile da una VRAM maggiore (circa 24 GB contro circa 12 GB in locale). Questo oggi impedisce di stimare con precisione la sensibilità a rumore, accenti e condizioni acustiche avverse. Impedisce anche un confronto controllato tra taglie di modello diverse.\cite{Radford2023Whisper}

Anche il modulo RAG mostra limiti strutturali. La qualità della risposta dipende dalla qualità dei chunk e dal rumore introdotto da OCR o ingestione non curata. Inoltre, anche con vincoli di grounding, possono restare generalizzazioni o risposte poco aderenti quando il materiale recuperato è incompleto o ambiguo.\cite{gao2024ragsurvey} La batteria tecnica sul backend rende questo problema più leggibile anche senza numeri aggregati. Il retrieval diretto può dare risultati buoni su query fattuali dirette e parafrasi semplici, ma l'esito conversazionale dipende ancora dall'intera catena: routing dell'intent, qualità del contesto recuperato e fase finale di generazione o riscrittura. Quando uno di questi passaggi si indebolisce, la risposta finale può diventare meno affidabile del retrieval sottostante.

Persistono infine limiti di portabilità pratica legati al runtime: su Windows la persistenza di ChromaDB può incorrere in lock e handle concorrenti, mitigati con procedure dedicate ma non eliminati del tutto. In generale, la riproducibilità dipende ancora dalla corretta installazione di dipendenze Python e dalla configurazione di componenti opzionali come OCR e servizi esterni.

\subsection{Sintesi rispetto alle domande di ricerca}
Le tre domande di ricerca, definite nel Capitolo~\ref{chapintroduzione}, possono essere rilette su tre assi: fattibilità tecnica, qualità percepita e riproducibilità. Per RQ1, il ciclo conversazionale integrato risulta fattibile sia nella build Windows nativa sia nella build WebGL, includendo memoria per-avatar e ingestione documentale con OCR. La lettura quantitativa resta parziale e va interpretata nel perimetro osservativo adottato.

Per RQ2, le evidenze qualitative suggeriscono che voice cloning, gestione del turn-taking (streaming + wait phrases) e segnali visivi di stato migliorano la fluidità percepita del dialogo. È ciò che emerge dal prototipo e dal comportamento osservato durante le sessioni analizzate.

Per RQ3, i risultati sono più netti: la scomposizione in micro-servizi, i contratti HTTP espliciti, gli health-check e la configurazione centralizzata degli endpoint supportano modularità e riproducibilità tra ambienti. Restano però limiti pratici, soprattutto sul runtime e sulle dipendenze, che impediscono ancora una riproducibilità pienamente ``push-button''.

La valutazione tecnica definisce ormai con una certa chiarezza il perimetro operativo del prototipo. Nella build Windows usata per le prove, la pipeline riesce a sostenere sessioni complete con una continuità sufficiente da giustificare un passaggio a una valutazione con utenti reali. Restano però limiti ben visibili: la latenza al primo audio continua a pesare sul turno, il retrieval risponde meglio a richieste specifiche che a domande vaghe, e in WebGL l'interazione diventa più fragile per l'effetto congiunto di rete, proxy e playback.

Da qui il passaggio successivo non riguarda più solo la tenuta della pipeline, ma il modo in cui il sistema viene percepito durante una sessione reale. Il capitolo seguente presenta quindi lo studio empirico comparativo svolto sul prototipo, mettendo a confronto una condizione neutra e una condizione personale.
\endgroup
